\documentclass[10pt,a4paper]{amsart}
\usepackage{amsmath,amssymb,amsthm,graphicx}
% Retained minimal preamble for the Stacks Project chapter more-algebra.tex (More on Algebra).
% The chapter file pinned-original-source.tex carries only \input{preamble} on its first line, so no
% preamble text lives inside the chapter itself.  The Stacks Project uses one single shared preamble for
% every chapter; the snapshot stages it identically as sources/stacks-chow/preamble.tex and
% sources/stacks-etale-cohomology/preamble.tex.  The block below is copied mechanically, line for line,
% from sources/stacks-chow/preamble.tex lines 14-18 (the Xy-pic setup) and 154-217 (the native theorem
% declarations together with the mathematical macros).  Nothing was added, removed or redefined; the
% wrapper class and the amsmath/amssymb/amsthm/graphicx base packages are supplied by the assembler.
% Chapter: more-algebra.tex -- More on Algebra.
% For commutative diagrams we use Xy-pic
\usepackage[all]{xy}

% We use 2cell for 2-commutative diagrams.
\xyoption{2cell}

% Theorem environments.
%
\theoremstyle{plain}
\newtheorem{theorem}[subsection]{Theorem}
\newtheorem{proposition}[subsection]{Proposition}
\newtheorem{lemma}[subsection]{Lemma}

\theoremstyle{definition}
\newtheorem{definition}[subsection]{Definition}
\newtheorem{example}[subsection]{Example}
\newtheorem{exercise}[subsection]{Exercise}
\newtheorem{situation}[subsection]{Situation}

\theoremstyle{remark}
\newtheorem{remark}[subsection]{Remark}
\newtheorem{remarks}[subsection]{Remarks}

\numberwithin{equation}{subsection}

% Macros
%
\def\lim{\mathop{\mathrm{lim}}\nolimits}
\def\colim{\mathop{\mathrm{colim}}\nolimits}
\def\Spec{\mathop{\mathrm{Spec}}}
\def\Hom{\mathop{\mathrm{Hom}}\nolimits}
\def\Ext{\mathop{\mathrm{Ext}}\nolimits}
\def\SheafHom{\mathop{\mathcal{H}\!\mathit{om}}\nolimits}
\def\SheafExt{\mathop{\mathcal{E}\!\mathit{xt}}\nolimits}
\def\Sch{\mathit{Sch}}
\def\Mor{\mathop{\mathrm{Mor}}\nolimits}
\def\Ob{\mathop{\mathrm{Ob}}\nolimits}
\def\Sh{\mathop{\mathit{Sh}}\nolimits}
\def\NL{\mathop{N\!L}\nolimits}
\def\CH{\mathop{\mathrm{CH}}\nolimits}
\def\proetale{{pro\text{-}\acute{e}tale}}
\def\etale{{\acute{e}tale}}
\def\QCoh{\mathit{QCoh}}
\def\Ker{\mathop{\mathrm{Ker}}}
\def\Im{\mathop{\mathrm{Im}}}
\def\Coker{\mathop{\mathrm{Coker}}}
\def\Coim{\mathop{\mathrm{Coim}}}

% Boxtimes
%
\DeclareMathSymbol{\boxtimes}{\mathbin}{AMSa}{"02}

%
% Macros for moduli stacks/spaces
%
\def\QCohstack{\mathcal{QC}\!\mathit{oh}}
\def\Cohstack{\mathcal{C}\!\mathit{oh}}
\def\Spacesstack{\mathcal{S}\!\mathit{paces}}
\def\Quotfunctor{\mathrm{Quot}}
\def\Hilbfunctor{\mathrm{Hilb}}
\def\Curvesstack{\mathcal{C}\!\mathit{urves}}
\def\Polarizedstack{\mathcal{P}\!\mathit{olarized}}
\def\Complexesstack{\mathcal{C}\!\mathit{omplexes}}
% \Pic is the operator that assigns to X its picard group, usage \Pic(X)
% \Picardstack_{X/B} denotes the Picard stack of X over B
% \Picardfunctor_{X/B} denotes the Picard functor of X over B
\def\Pic{\mathop{\mathrm{Pic}}\nolimits}
\def\Picardstack{\mathcal{P}\!\mathit{ic}}
\def\Picardfunctor{\mathrm{Pic}}
\def\Deformationcategory{\mathcal{D}\!\mathit{ef}}

% --- appended 2026-10-12 (Chief): the project own definitions of the reference/slogan/history
% display environments, copied mechanically from the shared Stacks preamble lines 8-12
% (\usepackage{verbatim} plus the three \newenvironment lines, which the chapter preamble keeps
% verbatim from the project preamble). Nothing added, removed or redefined: these environments
% carry no mathematical content and the project itself typesets them as comments. Without these
% three lines a counted statement whose own text uses them fails with Environment slogan undefined.
\usepackage{verbatim}
\newenvironment{reference}{\comment}{\endcomment}
\newenvironment{slogan}{\comment}{\endcomment}
\newenvironment{history}{\comment}{\endcomment}
\usepackage{cleveref}
\crefname{originallemma}{original lemma}{original lemmas}
\Crefname{originallemma}{Original lemma}{Original lemmas}
\makeatletter
\newlabel{lemma-more-relative-pseudo-coherent-is-moot}{{0E1T}{source}{}{sourcepointer-1}{}}
\newlabel{lemma-more-relative-pseudo-coherent-is-moot@cref}{{[originallemma][0][]0E1T}{[1][0][]source}{}{}{}}
\makeatother
\crefname{originallemma}{original lemma}{original lemmas}
\Crefname{originallemma}{Original lemma}{Original lemmas}
\makeatletter
\newlabel{lemma-perfect-over-regular-local-ring}{{09PC}{source}{}{sourcepointer-2}{}}
\newlabel{lemma-perfect-over-regular-local-ring@cref}{{[originallemma][0][]09PC}{[1][0][]source}{}{}{}}
\makeatother
\crefname{originallemma}{original lemma}{original lemmas}
\Crefname{originallemma}{Original lemma}{Original lemmas}
\makeatletter
\newlabel{lemma-flat-push-tor-amplitude}{{066J}{source}{}{sourcepointer-3}{}}
\newlabel{lemma-flat-push-tor-amplitude@cref}{{[originallemma][0][]066J}{[1][0][]source}{}{}{}}
\makeatother
\title{Tor amplitude over the polynomial presentation ring lifts from the fibre for pseudo-coherent complexes}
\author{The Stacks Project Authors}
\date{stacks-project:more-algebra.tex:local-snapshot-20261012; GNU FDL 1.2 or later}
\begin{document}
\maketitle
\setcounter{subsection}{1}
\noindent\textit{Editorial note.} All delivered mathematical text is the author's own, copied from the pinned Stacks Project chapter file more-algebra.tex; no statement or proof step was rewritten, reordered or removed. The counted span is the author environment and its own proof environment. References inside the span are handed to the assembler as navigation stubs that carry the canonical tag of the cited tag only; whether each of them is load-bearing is recorded on the stub itself.\par\medskip
\noindent\textit{Allowed standard prerequisites (editorial).}\begin{itemize}
\item A pseudo-coherent complex over A = R[x\_1, ..., x\_d]/I with A flat over R is pseudo-coherent as a complex of R[x\_1, ..., x\_d]-modules. invoked at proof lines 22193-22195. Reduces to the polynomial presentation ring.
\item For a map R\_p -> A\_q of local rings with A\_q a quotient of a regular local ring, a pseudo-coherent complex whose fibre has cohomology confined to [a, b] is perfect with tor amplitude in [a - d, b]. invoked at proof lines 22196-22199. Produces the perfectness and the shifted tor amplitude over A\_q.
\item A flat ring map pushes tor amplitude bounds through: a perfect complex over A\_q with tor amplitude in [a - d, b] over A\_q has tor amplitude in [a - d, b] over R\_p. invoked at proof lines 22200-22201. Transfers the tor amplitude down to R.
\end{itemize}\medskip
\ifdefined\hypertarget\hypertarget{sourcepointer-1}{}\fi
\noindent\textit{Source reference.} Original lemma 0E1T: sources/stacks-more-algebra/pinned-original-source.tex, source lines 21885--21885. Cited at proof lines 22193-22194 to reduce to the polynomial presentation; the conclusion is used inside the counted span.\par\medskip
\ifdefined\hypertarget\hypertarget{sourcepointer-2}{}\fi
\noindent\textit{Source reference.} Original lemma 09PC: sources/stacks-more-algebra/pinned-original-source.tex, source lines 20229--20229. Cited at proof lines 22196-22199 to produce perfectness with tor amplitude in [a - d, b] over A\_q; the conclusion is used inside the counted span.\par\medskip
\ifdefined\hypertarget\hypertarget{sourcepointer-3}{}\fi
\noindent\textit{Source reference.} Original lemma 066J: sources/stacks-more-algebra/pinned-original-source.tex, source lines 17100--17100. Cited at proof lines 22200-22201 to transfer the tor amplitude to R; the conclusion is used inside the counted span.\par\medskip
\begin{lemma}
\label{lemma-lift-from-fibre-relatively-perfect}
Let $R$ be a ring. Let $A = R[x_1, \ldots, x_d]/I$
be flat and of finite presentation over $R$.
Let $\mathfrak q \subset A$ be a prime ideal lying over
$\mathfrak p \subset R$. Let $K \in D(A)$ be pseudo-coherent.
Let $a, b \in \mathbf{Z}$. If
$H^i(K_\mathfrak q \otimes_{R_\mathfrak p}^\mathbf{L} \kappa(\mathfrak p))$
is nonzero only for $i \in [a, b]$, then
$K_\mathfrak q$ has tor amplitude in $[a - d, b]$ over $R$.
\end{lemma}

\begin{proof}
By Lemma \ref{lemma-more-relative-pseudo-coherent-is-moot}
$K$ is pseudo-coherent as a complex of $R[x_1, \ldots, x_d]$-modules.
Therefore we may assume $A = R[x_1, \ldots, x_d]$.
Applying Lemma \ref{lemma-perfect-over-regular-local-ring}
to $R_\mathfrak p \to A_\mathfrak q$ and the complex $K_\mathfrak q$
using our assumption, we find that $K_\mathfrak q$ is perfect
in $D(A_\mathfrak q)$ with tor amplitude in $[a - d, b]$.
Since $R_\mathfrak p \to A_\mathfrak q$ is flat, we conclude
by Lemma \ref{lemma-flat-push-tor-amplitude}.
\end{proof}

\end{document}
